%!TEX program = xelatex
\documentclass[12pt,a4paper]{ctexart}

\usepackage{geometry}
\geometry{left=2.5cm,right=2.5cm,top=2.5cm,bottom=2.5cm}
\usepackage{amsmath,amssymb,amsthm}
\usepackage{booktabs,tabularx,multirow}
\usepackage{float,enumitem,xcolor}
\usepackage{hyperref}

\theoremstyle{definition}
\newtheorem{definition}{定义}[section]
\newtheorem{theorem}{定理}[section]
\newtheorem{lemma}{引理}[section]
\newtheorem{proposition}{命题}[section]
\newtheorem{corollary}{推论}[section]

\newcommand{\R}{\mathbb{R}}
\newcommand{\Z}{\mathbb{Z}}
\newcommand{\set}[1]{\{#1\}}
\newcommand{\abs}[1]{|#1|}
\newcommand{\VEC}[1]{\boldsymbol{#1}}
\newcommand{\ind}[1]{\mathbf{1}_{[#1]}}

\title{\textbf{方案三：列生成方法\\（Dantzig-Wolfe 分解）}}
\author{数学建模备赛}
\date{2026年7月19日}

\begin{document}
\maketitle

% ==================== 问题分析 ====================
\section{问题分析}

在 $10\times 10$ 的网格海域中，船舶从起始岛 $B(1,5)$ 出发，须在 $30$ 天内抵达终点交付岛 $E(10,5)$。海域中有 $5$ 个特殊功能点：三个作业点 $W_1,W_2,W_3$（各有不同收益率和连续作业上限）和两个补给站 $S_1,S_2$。初始资源为燃油 $35$、淡水 $45$、食物 $30$、资金 $240$、目标物资 $100$，载重上限 $120$。优化目标为字典序 $\max Z$、次 $\max M$。

\textbf{集合划分视角与列生成动机}。在数学规划的视角下，本题可视为一条"恰好选择一个可行航程"的集合划分问题。若将每一条完整的航行方案（包含路径序列、工作天数分配、各补给站采购量）视为一个决策"列" $r\in\mathcal{R}$，则原始问题直接写为：
\begin{equation}\label{eq:orig}
  \max\ \sum_{r\in\mathcal{R}} Z(r)\cdot a_r,\quad \text{s.t.}\ \sum_{r\in\mathcal{R}} a_r = 1,\ a_r\in\{0,1\}
\end{equation}

问题 (\ref{eq:orig}) 的困难在于 $\mathcal{R}$（所有可行航程）的规模极大——路径骨架的组合数（$5^k$, $k\le 8$）乘以每个骨架的工作天数组合（指数级），无法显式枚举。然而，在最优解中只需恰好一个 $a_r=1$，其余的 $a_r$ 均为 0。列生成（Dantzig-Wolfe 分解）正是为这种情况设计的：仅在需要时动态生成能改进目标的列，通过主问题和定价子问题的交替迭代逐步逼近全局最优。

\textbf{本题列生成的特殊性}。通常列生成的优势在于定价子问题比原问题更容易求解（如最短路径问题）。但在本题中，定价子问题本身是 NP-hard 的带资源约束最短路问题（RCSP），求解策略（骨架枚举+贪婪模拟）与直接求解原始问题并无本质区别。因此本题中使用列生成的价值更多在于方法论——展示经典的 Dantzig-Wolfe 分解框架如何应用于一个具体的组合优化问题，以及为 $30\times 30$ 网格规模的任务 3 铺路（届时全枚举不可行，列生成的动态列生成将展现真正的规模优势）。

% ==================== 模型构建与论证 ====================
\section{模型构建与论证}

\subsection{主问题（Restricted Master Problem）}

设当前列池为 $\mathcal{R}'\subseteq\mathcal{R}$（初始仅含 $B\to E$ 直航列），受限主问题（RMP）为：
\begin{equation}\label{eq:rmp}
  \max\ \sum_{r\in\mathcal{R}'} Z(r)\cdot a_r,\quad \text{s.t.}\ \sum_{r\in\mathcal{R}'} a_r = 1,\ a_r\in[0,1]
\end{equation}

由于仅一条凸性约束（$\sum a_r=1$），RMP 的 LP 松弛最优解极其简单——在 $\mathcal{R}'$ 中选择 $Z$ 值最大的列，令其 $a_r=1$，其余为 0。对应的对偶变量 $\pi = \max_{r\in\mathcal{R}'} Z(r)$。

\begin{proposition}[主问题的退化性]\label{prop:degenerate}
  当主问题仅含一条凸性约束时，任意基矩阵为 $1\times 1$，对偶变量 $\pi$ 恰好等于当前池中最优列的 $Z$ 值，每次列更新严格提升 $\pi$。
\end{proposition}

\subsection{定价子问题（Pricing Subproblem）}

对任意列 $r\notin\mathcal{R}'$，其 LP 检验数（reduced cost）为：
\begin{equation}
  \bar{c}_r = Z(r) - \pi
\end{equation}

若 $\bar{c}_r > 0$，则该列可改进 RMP 的目标。定价子问题即为：
\begin{equation}\label{eq:pricing}
  \text{寻找可行航程 } r \text{ 使得 } Z(r) > \pi,\ \text{或证明不存在}
\end{equation}

\begin{theorem}[列生成的有限终止性]\label{thm:cg-conv}
  在路径骨架空间有限（542 种可行骨架，含 $0\sim 7$ 个中间点）、工作天数组合有限的条件下，列生成算法在有限次迭代后必然终止，终止时所得解等于全枚举最优解。
\end{theorem}

\begin{proof}
  每次定价迭代若找到改进列，则 $\pi_{\text{new}} = Z(r) > \pi_{\text{old}}$——$\pi$ 严格单调递增。$Z$ 的上界由全部 30 天 $\times$ 最高收益率（$28$/天）确定（约 $940$）。池的大小受限于骨架数 $\times$ 工作组合数（有限积）。因此算法必在有限步内终止。终止时 $\nexists r: Z(r) > \pi$，即 $\pi = \max_{r\in\mathcal{R}} Z(r)$。命题 \ref{prop:degenerate} 保证每次迭代的 $\pi$ 取自池中最优列，故池中最优列的 $Z$ 即全局最大 $Z$。\end{proof}

\subsection{定价子问题的求解}

定价子问题本身是 NP-hard 的 RCSP（带资源约束的最短路问题）。在本问题规模下，采用"骨架全枚举 + 工作天数枚举 + 贪婪前向模拟"精确求解：

\begin{enumerate}
  \item \textbf{外层}：生成所有路径骨架 $B\to\{W_1,W_2,W_3,S_1,S_2\}^*\to E$（过滤相邻重复、旅行超限）。
  \item \textbf{中层}：对路径中的 $n_w$ 个作业点访问，枚举总工作天数组合 $(W_1,\dots,W_{n_w})$。
  \item \textbf{内层}：对每个 $(P_i,W)$ 对调用 \texttt{gsim} 进行贪婪前向模拟——逐日消耗资源，补给日按"恰好支撑到下一补给站"采购，检验载重和非负约束。若 $Z>\pi$ 则输出该列。
\end{enumerate}

\textbf{含停泊策略的定价求解}。在最终版（\texttt{solve\_cg.m}）中，工作天数枚举的上界从 $WM_j$ 扩展至 $WM_j\times 3$（支持最多 3 个作业块）。日历天占用由分段策略计算：若 $W_j>WM_j$，分 $\lceil W_j/WM_j\rceil$ 段，段间各 1 天停泊。这使得定价子问题能发现含停泊-再作业的更优列。

\subsection{算法流程}

\begin{enumerate}
  \item 初始化列池 $\mathcal{R}' = \{B\to E\}$，$\pi = \max_{r\in\mathcal{R}'} Z(r)$。
  \item 定价迭代：枚举全部骨架和工作天数组合，若发现 $Z(r)>\pi$ 的可行航程，则 $r$ 加入 $\mathcal{R}'$，更新 $\pi$，标记 $\text{found}=\text{true}$。
  \item 若整个枚举空间均无 $Z(r)>\pi$ 的列，则 $\text{found}=\text{false}$，算法终止。
  \item 输出 $\mathcal{R}'$ 中 $Z$ 最大的列作为最优解。
\end{enumerate}

\subsection{对偶变量的经济学解释}

对偶变量 $\pi$ 在此特殊问题结构中具有直观的经济学含义：
\begin{itemize}
  \item $\pi$ 代表当前列池中"最佳航程的价值"（影子价格）。
  \item 检验数 $Z(r)-\pi$ 衡量一条新航程相对于当前最佳航程的增量价值——若为正，该航程"值得加入池中"。
  \item $\pi$ 的每次跃升对应一条质量更高的航程被发现，体现了从"直接驶向 E"（$Z=100$）到"逐步开发作业点"（$Z$ 逐步突破至 $328$）的搜索进程。
\end{itemize}

% ==================== 数值计算 ====================
\section{数值计算}

\subsection{实验设置}

在 MATLAB 环境下运行 \texttt{solve\_cg.m}（含停泊-再作业策略的最终版），骨架枚举空间为 $542$ 种可行骨架（含 $0\sim 7$ 个中间点）。

\subsection{实验结果}

\begin{table}[H]
  \centering
  \caption{列生成实验结果}
  \begin{tabular}{lc}
    \toprule
    指标 & 含停泊策略（最终版） \\
    \midrule
    最优 $Z$    & $\mathbf{328}$ \\
    最优 $M$    & $21$ \\
    路径        & $B\!\to\!W_1\!\to\!S_2\!\to\!W_3\!\to\!S_2\!\to\!E$ \\
    $W_1$ 工作  & $3$ 天 \\
    $W_3$ 工作  & $6$ 天（一次造访，$3+1_{\text{停泊}}+3$） \\
    旅行天数    & $19$ \\
    总天数      & $29$（1 天裕量） \\
    生成列数    & $12$ \\
    定价迭代    & $13$ \\
    \bottomrule
  \end{tabular}
\end{table}

\subsection{对偶变量演化}

对偶变量 $\pi$ 的演化从 $Z=100$（仅含直航列）开始，每次定价迭代发现改进列后 $\pi$ 跃升至新的 $Z$ 值，最终收敛至 $328$。每次跃迁代表一条质量更高的航程被加入池中。$\pi$ 的单调递增性由定价子问题的结构保证（每次找到的更优列必有 $Z>\pi$）。

\subsection{最优方案的逐日明细}

\begin{table}[H]
  \centering
  \caption{最优方案逐日明细（$Z=328$, $M=21$）}
  \begin{tabular}{cclc}
    \toprule
    日 & 位置 & 行为 & 备注 \\
    \midrule
    1--3   & $(1,5)\to(2,7)$  & 移动          & $B\to W_1$ \\
    4--6   & $(2,7)$          & $W_1$ 作业    & $+60$ Z \\
    7--12  & $(2,7)\to(7,6)$  & 移动          & $W_1\to S_2$（含 1 天到达 $S_2$）\\
    12     & $(7,6)$          & $S_2$ 补给    & 采购 O=41, H=37, F=28（花费 175，$M=65$）\\
    13--15 & $(7,6)\to(8,8)$  & 移动          & $S_2\to W_3$ \\
    16--18 & $(8,8)$          & $W_3$ 作业第 1 段 & $+84$ Z \\
    19     & $(8,8)$          & 停泊          & 重置连续计数器 \\
    20--22 & $(8,8)$          & $W_3$ 作业第 2 段 & $+84$ Z \\
    23--25 & $(8,8)\to(7,6)$  & 移动          & $W_3\to S_2$ \\
    25     & $(7,6)$          & $S_2$ 补给    & 采购 O=8, H=12, F=8（花费 44，$M=21$）\\
    26--29 & $(7,6)\to(10,5)$ & 移动          & $S_2\to E$ \\
    \bottomrule
  \end{tabular}
\end{table}

全程载重最大 $117\le 120$，资源恒非负，作业连续不超限（$W_3$：$3$ 天作业→1 天停泊重置→$3$ 天作业，每段均不超 $3$ 天上限）。

% ==================== 结果分析 ====================
\section{结果分析}

\subsection{停泊策略的价值量化}

含停泊策略相对于无停泊策略的提升：
\begin{itemize}
  \item $Z$：$288\to 328$（$+40$，$+13.9\%$）
  \item 旅行天数：$19$
  \item 总天数：$29$（有 1 天缓冲）
  \item 列数：$12$
  \item 迭代：$13$
\end{itemize}

停泊-再作业策略将 $W_3$ 的两次造访合并为一次，$W_1$ 得以从 1 天增至 3 天。

\subsection{收敛效率分析}

含停泊策略仅用 $12$ 个生成列和 $13$ 次迭代即收敛，反而少于无停泊策略的 $15/16$ 次。原因在于停泊策略拓展了可行域：更高质量的列（如 $Z=328$ 的列）无需经过中间阶梯即可被定价子问题发现，导致算法"跳过"了一些中间改进步骤。

\subsection{与方案一（全枚举）的关系}

方案一的搜索空间（$\sim 5\times 10^4$ 骨架 $\times\sim 200$ 工作组合 $\approx 1\times 10^7$ 次贪婪模拟）远大于列生成（$\sim 5\times 10^4$ 骨架 $\times\sim 200$ 组合，但仅在定价子问题中调用，且随 $\pi$ 上升而过滤能力增强）。在本题规模下列生成的 overhead（池管理和对偶更新）大于纯枚举。列生成在本题中的主要价值不是加速，而是：
\begin{enumerate}
  \item 提供对偶理论的最优性保证（定理 \ref{thm:cg-conv}），从另一角度验证结果。
  \item 为 $30\times 30$ 网格的任务 3 提供可扩展的框架——届时显式枚举 $5^k$ 不再可行，但列生成可通过设计更智能的定价算法（如标签设置算法）避免全枚举。
\end{enumerate}

\subsection{方法评价}

\textbf{优势}：
\begin{itemize}
  \item 对偶理论严格保证全局最优性（在有限骨架空间内）。
  \item 天然支持并行化——定价子问题可在不同骨架上并行搜索。
  \item 框架可扩展至 $30\times 30$ 网格。
  \item 对偶信息（$\pi$ 的演化）提供了问题结构的洞察。
\end{itemize}

\textbf{局限}：
\begin{itemize}
  \item 本题枚举空间小（$542$ 骨架），列生成的管理开销大于直接枚举。
  \item 定价子问题实质仍是枚举，未利用经典列生成中"子问题更易"的优势。
  \item 主问题仅一条约束导致退化——列生成的典型优势（处理大量约束）在本题中未体现。
\end{itemize}

\end{document}
